\subsection{Description of a Semisimple Module}

For the remainder of this section, A is a ring, and S is the set of classes of simple A-modules. For every \(\lambda \in S\) , we choose a simple module \(S_{\lambda}\) of class \(\lambda\) (for example, \(S_{\lambda} = \lambda\) ); we denote the opposite ring of the field of endomorphisms of \(S_{\lambda}\) by \(D_{\lambda}\) . We view \(S_{\lambda}\) as an \((A, D_{\lambda})\) -bimodule.

Let \(M\) be an A-module. For every \(\lambda \in \mathcal{S}\) , \(\mathrm{Hom}_{\mathrm{A}}(\mathrm{S}_{\lambda}, M)\) is a left vector space over the field \(D_{\lambda}\) . By VIII, p.~59, and II, §1, No.~6, p.~202, Proposition 6, there exists a unique A-linear mapping, called canonical,

\[
\alpha_ {\mathrm{M}}: \bigoplus_ {\lambda \in \mathcal {S}} \left(\mathrm{S} _ {\lambda} \otimes_ {\mathrm{D} _ {\lambda}} \operatorname{Hom} _ {\mathrm{A}} (\mathrm{S} _ {\lambda}, \mathrm{M})\right) \to \mathrm{M}
\]

satisfying the relation

\[
\alpha_ {\mathrm{M}} (s \otimes f) = f (s)\tag{10}
\]

for \(\lambda \in \mathcal{S}\) , \(s \in S_{\lambda}\) , and \(f \in \mathrm{Hom}_{\mathrm{A}}(\mathrm{S}_{\lambda}, \mathrm{M})\) . If we endow \(\bigoplus_{\lambda \in \mathcal{S}} (\mathrm{S}_{\lambda} \otimes_{\mathrm{D}_{\lambda}} \mathrm{Hom}_{\mathrm{A}}(\mathrm{S}_{\lambda}, \mathrm{M}))\) and \(\mathrm{M}\) with their natural structures of \(\mathrm{End}_{\mathrm{A}}(\mathrm{M})\) -modules, then the mapping \(\alpha_{\mathrm{M}}\) is \(\mathrm{End}_{\mathrm{A}}(\mathrm{M})\) -linear.

PROPOSITION 7. --- Let M be an A-module. The canonical mapping \(\alpha_{\mathrm{M}}\) is injective. For every \(\lambda \in \mathcal{S}\) , the image under \(\alpha_{\mathrm{M}}\) of \(\mathrm{S}_{\lambda} \otimes_{\mathrm{D}_{\lambda}} \mathrm{Hom}_{\mathrm{A}}(\mathrm{S}_{\lambda}, \mathrm{M})\) is the isotypical component of M of type \(\lambda\) . The image of \(\alpha_{\mathrm{M}}\) is the socle of M. The A-module M is semisimple if and only if the mapping \(\alpha_{\mathrm{M}}\) is bijective.

Let \(\lambda \in \mathcal{S}\) . Denote the isotypical component of M of type \(\lambda\) by \(\mathrm{M}_{\lambda}\) . Every A-linear mapping from \(\mathrm{S}_{\lambda}\) to M takes its values in \(\mathrm{M}_{\lambda}\) (VIII, p.~66, Proposition 5). Consequently, by Proposition 3, a) of VIII, p.~62, the mapping \(\alpha_{\mathrm{M}}\) induces a bijection from \(\mathrm{S}_{\lambda} \otimes_{\mathrm{D}_{\lambda}} \mathrm{Hom}_{\mathrm{A}}(\mathrm{S}_{\lambda}, \mathrm{M})\) to \(\mathrm{M}_{\lambda}\) . The proposition follows because the socle of M is the direct sum of the family \((\mathrm{M}_{\lambda})_{\lambda \in \mathcal{S}}\) and the module M is semisimple if and only if it is equal to its socle.

DEFINITION 5. --- Let M be a semisimple A-module. A description of M (with respect to the family \((\mathrm{S}_{\lambda})_{\lambda\in\mathcal{S}}\) ) is a pair \(((\mathrm{V}_{\lambda})_{\lambda\in\mathcal{S}},\alpha)\) , where \(V_{\lambda}\) is a left vector space over the field \(D_{\lambda}\) for each \(\lambda\in S\) and \(\alpha:\bigoplus_{\lambda\in\mathcal{S}}(\mathrm{S}_{\lambda}\otimes_{\mathrm{D}_{\lambda}}\mathrm{V}_{\lambda})\to\mathrm{M}\) is an isomorphism of A-modules.

By Proposition 7, every semisimple A-module M has a canonical description: it is the pair \(\left(\left(\mathrm{Hom}_{\mathrm{A}}\left(\mathrm{S}_{\lambda},\mathrm{M}\right)\right)_{\lambda \in \mathcal{S}},\alpha_{\mathrm{M}}\right)\) , where \(\alpha_{\mathrm{M}}\) is the A-linear mapping defined by formula (10).

PROPOSITION 8. --- Let M be a semisimple A-module and \(((\mathrm{V}_{\lambda})_{\lambda \in \mathcal{S}}, \alpha)\) a description of M.

\begin{enumerate}
\def\labelenumi{\alph{enumi})}
\item
  For every \(\lambda \in \mathcal{S}\) , the mapping \(\alpha\) induces an isomorphism from the A-module \(S_{\lambda} \otimes_{D_{\lambda}} V_{\lambda}\) to the isotypical component of M of type \(\lambda\) .
\item
  For every \(\lambda \in \mathcal{S}\) , the mapping \(\beta_{\lambda} : \mathrm{V}_{\lambda} \to \mathrm{Hom}_{\mathrm{A}}(\mathrm{S}_{\lambda}, \mathrm{M})\) defined by \(\beta_{\lambda}(v)(s) = \alpha(s \otimes v)\) is \(\mathrm{D}_{\lambda}\) -linear and bijective.
\item
  Let \(\mathbf{N}\) be a submodule of \(\mathbf{M}\) . There exists a unique family \((\mathrm{W}_{\lambda})_{\lambda \in \mathcal{S}}\) with the following properties: \(\mathrm{W}_{\lambda}\) is a \(\mathrm{D}_{\lambda}\) -linear subspace of \(\mathrm{V}_{\lambda}\) for every \(\lambda \in \mathcal{S}\) , and \(\mathbf{N}\) is the image under \(\alpha\) of the module \(\bigoplus_{\lambda \in \mathcal{S}} (\mathrm{S}_{\lambda} \otimes_{\mathrm{D}_{\lambda}} \mathrm{W}_{\lambda})\) identified with its canonical image in the module \(\bigoplus_{\lambda \in \mathcal{S}} (\mathrm{S}_{\lambda} \otimes_{\mathrm{D}_{\lambda}} \mathrm{V}_{\lambda})\) . For every \(\lambda \in \mathcal{S}\) , \(\mathrm{W}_{\lambda}\) is the set of elements \(v \in \mathrm{V}_{\lambda}\) such that \(\alpha(s \otimes v)\) belongs to \(\mathbf{N}\) for every \(s \in \mathrm{S}_{\lambda}\) .
\end{enumerate}

The A-module \(S_{\lambda} \otimes_{D_{\lambda}} V_{\lambda}\) is isotypical of type \(\lambda\) for every \(\lambda \in S\) . Assertion a) then follows from the facts that \(\alpha\) is an isomorphism and that M is the direct sum of the family \((M_{\lambda})_{\lambda \in \mathcal{S}}\) (VIII, p.~65, Proposition 4, a)).

Let \(\lambda \in \mathcal{S}\) . Denote the restriction of \(\alpha\) to \(S_{\lambda} \otimes_{D_{\lambda}} V_{\lambda}\) by \(\alpha_{\lambda}: S_{\lambda} \otimes_{D_{\lambda}} V_{\lambda} \to M\) . With the notation of No.~3 applied to the \((A, D_{\lambda})\) -bimodule \(S_{\lambda}\) , we have

\[
\beta_ {\lambda} = \gamma (\alpha_ {\lambda}) = \mathcal {H} (\alpha_ {\lambda}) \circ \beta_ {\mathrm{V} _ {\lambda}},
\]

where the last equality follows from VIII, p.~60. Hence, \(\beta_{\lambda}\) is the composition of the \(D_{\lambda}\) -linear homomorphism \(\mathcal{H}(\alpha_{\lambda})\) and \(\beta_{V_{\lambda}}\) . Assertion b) then follows from Proposition 3, b) of VIII, p.~62.

Let \(N\) be a submodule of \(M\) . We have \(N = \bigoplus_{\lambda \in \mathcal{S}} (N \cap M_{\lambda})\) (VIII, p.~66, Corollary), so c) follows from Theorem 2 of VIII, p.~62.

COROLLARY. --- Let M be a semisimple A-module. For every submodule N of M and every element \(\lambda\) of S, identify \(\operatorname{Hom}_{\mathrm{A}}(\mathrm{S}_{\lambda},\mathrm{N})\) with the \(D_{\lambda}\) -linear subspace of \(\operatorname{Hom}_{\mathrm{A}}(\mathrm{S}_{\lambda},\mathrm{M})\) consisting of the mappings with image contained in N.

\begin{enumerate}
\def\labelenumi{\alph{enumi})}
\item
  The mapping \(\mathrm{N}\mapsto(\mathrm{Hom}_{\mathrm{A}}(\mathrm{S}_{\lambda},\mathrm{N}))_{\lambda\in\mathcal{S}}\) is a bijection from the set of A-submodules of M to the set of families \((\mathrm{W}_{\lambda})_{\lambda\in\mathcal{S}}\) such that for every \(\lambda\in S\) , \(W_{\lambda}\) is a \(D_{\lambda}\) -linear subspace of \(\mathrm{Hom}_{\mathrm{A}}(\mathrm{S}_{\lambda},\mathrm{M})\) .
\item
  The inverse bijection sends a family \((\mathrm{W}_{\lambda})_{\lambda \in \mathcal{S}}\) to the A-submodule \(\sum_{\lambda \in \mathcal{S}}\sum_{w\in \mathrm{W}_{\lambda}}w(\mathrm{S}_{\lambda})\) of M.
\end{enumerate}

This is a reformulation of Proposition 8, c) applied to the canonical description of M.

PROPOSITION 9. --- Let M and M' be semisimple A-modules, and consider descriptions \(((\mathrm{V}_{\lambda})_{\lambda\in\mathcal{S}},\alpha)\) and \(((\mathrm{V}_{\lambda}^{\prime})_{\lambda\in\mathcal{S}},\alpha^{\prime})\) of M and M', respectively. For every family \(\boldsymbol{f}=(f_{\lambda})_{\lambda\in\mathcal{S}}\) in \(\prod_{\lambda\in\mathcal{S}}\mathrm{Hom}_{\mathrm{D}_{\lambda}}(\mathrm{V}_{\lambda},\mathrm{V}_{\lambda}^{\prime})\) , there exists a unique A-linear mapping \(\varphi(\boldsymbol{f})\in\mathrm{Hom}_{\mathrm{A}}(\mathrm{M},\mathrm{M}^{\prime})\) for which the following diagram commutes:

\[
\begin{array}{c} \bigoplus_ {\lambda \in \mathscr {S}} (S _ {\lambda} \otimes_ {D _ {\lambda}} V _ {\lambda}) \xrightarrow {\alpha} M \\ \bigoplus_ {\lambda \in \mathscr {S}} (1 _ {S _ {\lambda}} \otimes f _ {\lambda}) \Bigg | _ {\downarrow} \qquad \qquad \qquad \qquad \qquad \qquad \Bigg | _ {\varphi (\boldsymbol {f})} \\ \bigoplus_ {\lambda \in \mathscr {S}} (S _ {\lambda} \otimes_ {D _ {\lambda}} V _ {\lambda} ^ {\prime}) \xrightarrow {\alpha^ {\prime}} M ^ {\prime}. \end{array}
\]

The mapping \(\varphi : \prod_{\lambda \in \mathcal{S}} \mathrm{Hom}_{\mathrm{D}_{\lambda}}(\mathrm{V}_{\lambda}, \mathrm{V}_{\lambda}') \to \mathrm{Hom}_{\mathrm{A}}(\mathrm{M}, \mathrm{M}')\) defined this way is a group isomorphism. When we have \(\mathrm{M} = \mathrm{M}'\) , \(\mathrm{V}_{\lambda} = \mathrm{V}_{\lambda}'\) for every \(\lambda \in \mathcal{S}\) , and \(\alpha = \alpha'\) , the mapping \(\varphi\) is a ring isomorphism from \(\prod_{\lambda \in \mathcal{S}} \mathrm{End}_{\mathrm{D}_{\lambda}}(\mathrm{V}_{\lambda})\) to \(\mathrm{End}_{\mathrm{A}}(\mathrm{M})\) .

In view of the description of the isotypical components of M and M' given in Proposition 8, a), this follows from Theorem 3 of VIII, p.~64 and Proposition 5, b) of VIII, p.~66.

COROLLARY. --- Let M be a semisimple A-module and M' an A-module. The mapping \(u \mapsto (\operatorname{Hom}(1_{\mathrm{S}_{\lambda}}, u))_{\lambda \in \mathcal{S}}\) from \(\operatorname{Hom}_{\mathrm{A}}(\mathrm{M}, \mathrm{M}')\) to

\[
\prod_ {\lambda \in \mathcal {S}} \mathrm{Hom} _ {\mathrm{D} _ {\lambda}} (\mathrm{Hom} _ {\mathrm{A}} (\mathrm{S} _ {\lambda}, \mathrm{M}), \mathrm{Hom} _ {\mathrm{A}} (\mathrm{S} _ {\lambda}, \mathrm{M} ^ {\prime}))
\]

is a group isomorphism. When \(\mathrm{M}'\) is equal to \(\mathrm{M}\) , it is an isomorphism from the ring \(\operatorname{End}_{\mathrm{A}}(\mathrm{M})\) to the ring \(\prod_{\lambda \in \mathcal{S}} \operatorname{End}_{\mathrm{D}_{\lambda}}(\operatorname{Hom}_{\mathrm{A}}(\mathrm{S}_{\lambda}, \mathrm{M}))\) .

This is a reformulation of Proposition 9 applied to the canonical descriptions of M and of the socle of M'.

